\chapter{Tumor Marker Tests}
\section*{What Is This Test?} Tumor marker tests measure substances or genetic features associated with some cancers in blood, urine, tissue, or other fluids.\section*{Alternative Names} Cancer biomarker testing; tumor marker blood test; oncologic biomarker.\section*{Principle} Immunoassay, molecular, or pathology methods detect a marker produced by tumor cells or the body's response.\section*{Why This Test Is Ordered} Markers may support diagnosis, prognosis, treatment selection, or monitoring when used with examination and imaging.\section*{Primary vs Secondary Test} Most are secondary tests; many are not accurate enough to screen asymptomatic people or diagnose cancer alone.\section*{How This Test Is Performed} Sample type and method vary by marker; tissue biomarkers may require biopsy.\section*{What Preparation Is Needed} Test-specific instructions apply. Report cancer history, treatments, pregnancy, and medicines.\section*{Understanding Results} Benign disease can raise markers and cancer can occur with normal values; trends and assay consistency are often more useful than one value.\section*{Follow-up Tests} Imaging, biopsy, pathology biomarkers, repeat marker testing, or oncology review may follow.\section*{References} \begin{enumerate}\item MedlinePlus. Tumor Marker Tests.\item National Cancer Institute. Tumor markers fact sheet.\end{enumerate}\clearpage\section*{Understanding Results and Follow-up}
The laboratory report should identify the specimen, method, units, reference interval or decision limit, and whether the result is positive, negative, abnormal, or inconclusive. Reference intervals vary with age, sex, pregnancy, specimen type, assay, and laboratory; a result outside a range is not automatically a diagnosis, and a result within a range does not always exclude disease. Discuss unexpected results with the ordering clinician, who can decide whether repeat or confirmatory testing, treatment, or referral is appropriate. Do not change medicines or delay urgent care based on an isolated result.
\section*{Regulatory Status and Authorization Date}This chapter describes a test category, not a specific commercial product. FDA, CDSCO, EU IVDR, and MHRA approval or registration dates therefore cannot be assigned without the assay manufacturer, product name, instrument, and intended use. For laboratory-developed tests, an individual agency approval date may not exist. See the Preface for the interpretation of regulatory status.
\clearpage
